% The text of the unresolved-reference fixtures. Every label and every cited key is defined, and the plain
% question marks of the second paragraph, the listings, the table, and the figure must never count as marks.
\section{Introduction}\label{sec:intro}
Flaky tests pass and fail on the same code~\cite{luo2014}. Detectors rerun each failure~\cite{luo2014,bell2018},
and a survey covers the field~\cite{parry2021}. Section~\ref{sec:method} describes the method,
Fig.~\ref{fig:overview} gives an overview, Table~\ref{tab:results} lists the results, and
Eq.~(\ref{eq:score}) defines the score, also written as \eqref{eq:score}. Algorithm~\ref{alg:rank} ranks the
failures, RQ\ref{rq:accuracy} asks how accurate the ranking is, and \autoref{sec:discussion} and
\cref{sec:method} close the paper on page~\pageref{sec:discussion}.

A developer asked, ``Why?? The test passed yesterday!'' Another wrote ``Really?!'' and a third asked, ``Is this
the right page??'' The value $f(?)$ stays open, and the most common (?) cause is the order of the tests. Is the
cause known? Not yet. The ternary form \texttt{a ? b : c} and the query \texttt{?q=flaky} appear in the logs.

\begin{enumerate}
\item\label{rq:accuracy} \textbf{RQ1:} How accurate is the ranking? (\textbf{RQ1})
\end{enumerate}
\begin{equation}\label{eq:score}
s(t) = \frac{f(t)}{r(t)} \quad (?)
\end{equation}

\section{Method}\label{sec:method}
\textbf{What now??} The tool ranks the failures. \textbf{Is it flaky?} The history decides.
\begin{verbatim}
grep -E "[?]" log.txt   # a literal ?? or (?) in the log
\end{verbatim}
\begin{lstlisting}[language={[Sharp]C},basicstyle=\ttfamily\small,keywordstyle=\bfseries]
var name = user?.Name ?? "unknown";
\end{lstlisting}
\begin{table}[h]\caption{Results}\label{tab:results}\centering
\begin{tabular}{lrr}\hline \textbf{Project} & \textbf{Accuracy} & \textbf{Flaky?}\\\hline
A & 0.71 & yes\\ B & ? & (?)\\\hline\end{tabular}
\end{table}
\begin{figure}[h]\centering
\begin{tikzpicture}\node[draw,diamond] {\textbf{?}};\end{tikzpicture}
\caption{Overview}\label{fig:overview}
\end{figure}
\begin{algorithm}[h]\caption{Rank}\label{alg:rank}
\begin{algorithmic}\STATE rank the failures\end{algorithmic}
\end{algorithm}

\section{Discussion}\label{sec:discussion}
Which failures need attention first? The ranking answers this question for 71\% of the builds.
